\section{\textit{VirtIO}}

O \textit{VirtIO} \cite{RustyRussellVirtIO} é um modelo de dispositivo virtual projetado para padronizar a comunicação de operações de \gls{io} entre \glspl{so} convidados e seus respectivos hipervisores.
Seu funcionamento assemelha-se ao de dispositivos físicos baseados em \gls{dma}, empregando regiões de memória compartilhada para a transferência direta de dados.

Para viabilizar essa comunicação, o modelo utiliza as chamadas \textit{Virtqueues} — filas circulares compartilhadas responsáveis pelo envio e recebimento das requisições. Dependendo das operações suportadas,
cada dispositivo virtual pode alocar uma ou mais dessas filas. Internamente, uma \textit{Virtqueue} divide-se em três componentes principais: uma tabela de descritores (\textit{Descriptor Table}), um anel de
índices disponíveis (\textit{Available Ring}) e um anel de índices já utilizados (\textit{Used Ring}).

Na prática, o ciclo de vida de uma operação de \gls{io} ocorre em duas etapas complementares. Inicialmente, quando o sistema convidado deseja realizar uma operação, o \textit{driver VirtIO} escreve os dados
na memória compartilhada, insere o descritor correspondente no \textit{Available Ring} e notifica o hipervisor sobre a nova pendência. 

Em resposta, o hipervisor consome os descritores desse anel, processa a operação solicitada e registra o resultado no \textit{Used Ring}. Por fim, uma interrupção é enviada ao sistema convidado para sinalizar a
conclusão da tarefa. Graças a esse design, múltiplas requisições podem ser tratadas de maneira assíncrona e em lote, o que minimiza as trocas de contexto e reduz expressivamente o custo inerente à virtualização.

Desse modo, o \textit{VirtIO} consolida-se como uma solução fundamental para a infraestrutura de virtualização moderna, pois alcança um equilíbrio entre a flexibilidade de uma interface padronizada e o
desempenho da paravirtualização.
